\section{Retained quadratic inventory validation}\label{sec:inventory}
For historical validation of the dynamic certificate, retain the quadratic inventory-adjustment economy below. Its numerical gain problem reduces to two scalar Riccati equations; the genuinely vector-state nonlinear experiment is Section~\ref{sec:nonlinearinventory}. With $M=\mathbf1\mathbf1'/d$,
\begin{equation}
 dX_s=(AX_s+a_s)ds+.1dW_s,\qquad A=-.2I+.3M,\label{eq:inventorydynamics}
\end{equation}
and minimize
\begin{equation}
 \mathcal J(t,x;a)=\E\!\left[\int_t^1\{X_s'(I+.5M)X_s+|a_s|^2\}ds+.5|X_1|^2\right].
 \label{eq:inventorycriterion}
\end{equation}
Controls are square-integrable adapted adjustments. Mean inventory affects drift and aggregate holding cost; idiosyncratic inventory also has its own adjustment dynamics and shocks. This is a coupled $d$-state dynamic problem, not an action maximization at frozen derivatives. It is a structured, uncalibrated application, and the structure is available to classical solvers as well.

The orthogonal and mean modes have $(a_\perp,q_\perp)=(-.2,1)$ and $(a_m,q_m)=(.1,1.5)$. Their Riccati gains satisfy
\[
 (p_i^*)'+2a_i p_i^*-(p_i^*)^2+q_i=0,\qquad p_i^*(1)=.5.
\]
We train two neural gain outputs $\widehat p_i(t)=.5+(1-t)\operatorname{softplus}(N_i(t))$ against these \emph{differential residuals}, not against exact Riccati solutions. The feedback is
\[
 \widehat a(t,x)=-\widehat P(t)x,\qquad
 \widehat P=\widehat p_\perp(I-M)+\widehat p_mM.
\]
The network changes its gain with time and controls all $d$ state coordinates. One learned pair is reused across the six dimensions; this reuse is not counted as six independent training runs.

\begin{theorem}[State-global quadratic neural policy certificate]\label{thm:inventory}
Suppose interval differentiation on a complete time cover bounds the two neural Riccati residuals by $\varepsilon_i$. Put $\kappa_\perp=1.4$, $\kappa_m=.8$, and
\[
 D=\max_i\varepsilon_i\frac{1-e^{-\kappa_i}}{\kappa_i},\qquad
 C=\frac{1-e^{-.8}}{.8}.
\]
For every $d\geq1$ and every $|x|^2/d\leq1$,
\begin{equation}
 0\leq\frac{\mathcal J(0,x;\widehat a)-\mathcal V(0,x)}d
 \leq D^2\left\{C+\frac{.01}{.8}(1-C)\right\}.
 \label{eq:inventorycertificate}
\end{equation}
The bound concerns the continuous-time controlled diffusion on all of $\R^d$, not a finite collection of initial states.
\end{theorem}
The proof uses a quadratic performance-difference identity and a closed-loop Lyapunov moment bound. It replaces state covering by algebra valid on the entire state space. The certificate's time-enclosure work depends on network size and the number of time cells, rather than a tensor grid in $d$. Applying the feedback requires a mean and coordinatewise operations, hence $O(d)$ work. This advantage comes from the model's two invariant modes; it is not a generic escape from dimensional dependence for nonlinear Bellman equations.

Ten fixed seeds, two widths, and both predeclared checkpoints are retained. Every final checkpoint has per-coordinate certified loss at most \RInventoryFinalMax; the complete range and costs appear in Table~\ref{tab:inventory}. The certificate scales to total loss by multiplication by $d$. The supplement also reports the earlier checkpoint and fixed-network time-cover refinements separately. A structure-aware classical calculation solves only two scalar Riccati equations. Its low cost is reported rather than hidden: this application establishes accurate, reproducibly generated neural feedback with full dynamic certification, not superiority over the classical solution of this specially structured problem.
\begin{table}[htbp]\centering\small
\caption{Certified neural policy loss per inventory coordinate}\label{tab:inventory}
\begin{tabular}{rrrrrr}
\toprule
Seed & Width & 250 updates & 1,000 updates & Train (s) & Verify (s)\\\midrule
16200 & 8 & 0.00117290 & 0.00009561 & 5.28 & 1.18\\
16201 & 16 & 0.00045470 & 0.00007661 & 7.05 & 2.85\\
16202 & 8 & 0.00051798 & 0.00006084 & 5.40 & 1.27\\
16203 & 16 & 0.00011929 & 0.00003574 & 7.44 & 3.26\\
16204 & 8 & 0.00179439 & 0.00008973 & 5.09 & 1.09\\
16205 & 16 & 0.00054105 & 0.00002530 & 6.28 & 3.05\\
16206 & 8 & 0.00036825 & 0.00004136 & 5.59 & 1.36\\
16207 & 16 & 0.00023854 & 0.00009053 & 6.34 & 3.34\\
16208 & 8 & 0.00099196 & 0.00005105 & 5.88 & 1.60\\
16209 & 16 & 0.00036174 & 0.00007427 & 5.08 & 2.29\\
\bottomrule\end{tabular}
\par\smallskip\footnotesize Each checkpoint uses a complete 512-cell time cover and a state-global quadratic proof. Bounds hold for every $|x_0|^2/d\leq1$. Multiply the reported bound by $d$ for total loss at $d=4,8,16,32,64,128$. One gain pair is reused across dimensions. Timings cover one training trajectory and its two checks, not six duplicated runs.
\end{table}

